% !TeX root = ../../Book3.tex
\section{调和函数}
\label{b3:sec:1501}

全纯函数的实部不只是满足一条微分方程。它在每个圆盘中心都等于边界的平均值；反过来，这种平均性质足以恢复光滑性和方程。先把这两个描述联系起来，再从正性得到定量控制。

\subsection{Laplace 方程}
\label{b3:subsec:150101}

设 \(D\subseteq\mathbb C=\mathbb R^2\) 为开集。实值函数 \(u\in C^2(D)\) 若满足 \(\Delta u=u_{xx}+u_{yy}=0\)，称为\term[diaohe-hanshu]{调和函数}[harmonic function]。复值函数调和是指实部、虚部均调和；最大值和正性结论以下只用于实值函数。\(x\)、\(x^2-y^2\)、\(\log|z|\) 分别在全平面、全平面和去心平面上调和。最后一个例子提醒我们：定义域的缺口可能承载奇性，不能在原点直接补值。

由 Cauchy--Riemann 方程和全纯函数的光滑性，\(f=u+\mathrm{i}v\) 全纯时 \(u_{xx}+u_{yy}=v_{yx}-v_{xy}=0\)。实变量的方程在二维保留了复变量的局部结构，但尚未保证全局调和共轭存在。

\subsection{平均值性质}
\label{b3:subsec:150102}

\term[pingjunzhi-xingzhi]{平均值性质}[mean-value property]把方程变成与尺度有关的等式。
\begin{theorem}[调和函数的平均值性质]\label{b3:thm:15-mean-value}
若 \(u\) 在 \(D\) 调和，且 \(\overline{B(a,r)}\subset D\)，则
\[
u(a)=\frac1{2\pi}\int_0^{2\pi}u(a+r e^{\mathrm{i}t})\,\mathrm{d}t
=\frac1{\pi r^2}\int_{B(a,r)}u\,\mathrm{d}A.
\]
反之，连续函数在每个内部圆盘上满足圆周平均等式，则为调和函数。
\end{theorem}
\begin{proof}
圆周平均记作 \(m(r)\)。微分和散度公式给出
\[
 m'(r)=\frac1{2\pi r}\int_{\partial B(a,r)}\partial_nu\,\mathrm{d}s
 =\frac1{2\pi r}\int_{B(a,r)}\Delta u\,\mathrm{d}A=0.
\]
由连续性 \(m(r)\to u(a)\)，得第一式；乘 \(2r\) 并积分得面积版本。
反向证明不能预先假定函数可微。在相对紧的小开集上取光滑径向非负核 \(\eta_\varepsilon\)，积分为一。极坐标和假设给 \(u*\eta_\varepsilon=u\)，故 \(u\) 局部光滑。对光滑函数在中心作二阶 Taylor 展开，圆周平均为 \(u(a)+r^2\Delta u(a)/4+o(r^2)\)。假设使二阶系数为零，故 \(\Delta u=0\)。
\end{proof}
径向磨光同时说明平均值性质为何能恢复所有阶导数：不是从连续性直接求导，而是先将函数识别成它的光滑卷积。

\subsection{最大值原理}
\label{b3:subsec:150103}

\begin{theorem}[强最大值原理]\label{b3:thm:15-maximum}
区域 \(D\) 上的实值调和函数若在内部取得全域最大值或最小值，则为常数。若 \(D\) 有界、\(u\in C(\overline D)\) 且在 \(D\) 调和，则最大值和最小值均可在边界取得。
\end{theorem}
\begin{proof}
设最大值为 \(M=u(a)\)。每个以 \(a\) 为中心的内部圆周上都有 \(u\le M\)，而平均等于 \(M\)。连续非负函数 \(M-u\) 的积分为零，故它在圆周恒零。所有充分小的半径合起来说明 \(u=M\) 的集合开；它由连续性又是闭集。连通性使其等于 \(D\)。最小值情形应用于 \(-u\)。有界域闭包紧，最大值存在；若极值只在内部出现，则上面的结论使函数处处相等，边界亦取得该值。
\end{proof}
有界性与边界条件各有作用。上半平面上的 \(u(x+\mathrm{i}y)=y\) 在实轴为零，却不恒零；没有无穷远控制，边界最大值原理和 Dirichlet 解的唯一性都会失败。

\subsection{Harnack 不等式}
\label{b3:subsec:150104}

为了比较不同点的正调和函数，先建立圆盘表示。在 \(|z|<R\) 记
\[
P_R(z,t)=\frac{R^2-|z|^2}{|Re^{\mathrm{i}t}-z|^2}.
\]
它是 \((Re^{\mathrm{i}t}+z)/(Re^{\mathrm{i}t}-z)\) 的实部，故关于 \(z\) 调和；几何级数表明其圆周平均为一。若 \(\varphi\) 连续，积分 \((2\pi)^{-1}\int P_R(z,t)\varphi(Re^{\mathrm{i}t})\,\mathrm{d}t\) 在内部调和。靠近边界一点时，远离该点的核一致趋零，近处由 \(\varphi\) 的一致连续性控制，所以积分连续取边值 \(\varphi\)。与 \(u\) 比较并用最大值原理，便得到 \(u\) 的 Poisson 表示。

\begin{theorem}[Harnack]\label{b3:thm:15-harnack}
若 \(u\ge0\) 在 \(B(0,R)\) 调和，则 \(|z|=r<R\) 时
\[
\frac{R-r}{R+r}\,u(0)\le u(z)\le\frac{R+r}{R-r}\,u(0).
\]
\end{theorem}
\begin{proof}
先在 \(r<S<R\) 的圆盘用上述表示。核满足 \((S-r)/(S+r)\le P_S(z,t)\le(S+r)/(S-r)\)，对非负边值积分后利用平均值性质得到不等式，再令 \(S\uparrow R\)。
\end{proof}
有限条相交圆盘构成的链把估计传到区域中任意固定紧集。这给出 Harnack 收敛原理：递增调和函数列若在一点有界，就局部一致收敛到调和函数；否则处处趋于正无穷。证明时把差函数代入链式估计，中心处的 Cauchy 性即变成紧集上的一致 Cauchy 性，再用平均值性质的逆命题。


平均值公式还给出内部导数估计。取 \(\overline{B(a,2r)}\subset D\)，在 \(B(a,r)\) 上用半径 \(2r\) 的 Poisson 表示。对核关于 \(z\) 的各阶导数求界，分母距离至少 \(r\)，得到
\[
\sup_{B(a,r)}|\partial^\alpha u|\le C_\alpha r^{-|\alpha|}\sup_{B(a,2r)}|u|.
\]
常数只依赖多重指标，因为把变量缩放到单位圆盘之后，核导数在半径二分之一的闭圆盘一致有界。估计同时控制导数与局部尺度；区域靠近边界时可用的 \(r\) 变小，不能把同一个内部常数当作边界导数界。

\ProseParagraphBreak

若调和函数族在每个紧集上统一有界，以上估计给导数的统一界和等度连续性。对一列这样的函数，先在第一个紧集抽出一致收敛子列，再逐层抽取并取对角列；所得局部一致极限保留平均等式，因此调和。对更高导数重用同一估计，可得到局部各阶导数收敛。这里的紧性机制与全纯正规族相似，但使用的是实值核和平均，而不是复曲线积分。

\ProseParagraphBreak

具体计算可以检验这些结论。边值 \(\varphi(t)=1+\cos t\) 的单位圆盘延拓为 \(1+r\cos\theta\)，非负，且在边界点 \(-1\) 为零。内部仍严格正，因为正调和函数只有在处处零时才可在内部为零。相反，\(u=x\) 在直径上有很多内部零点，却不违背强最小值原理，因为它在圆盘中不非负。不能把“调和函数无内部最小值”误读成“调和函数不能有内部零点”。

\ProseParagraphBreak

Harnack 原理中一点有界的量词也值得区分。若递增列 \(u_n\) 全部调和但不非负，可以先减 \(u_1\)，使差列非负；在一点上界便由 Harnack 链传到任何固定紧集。若没有单调性，一点有界不足以控制全列，例如 \(u_n(z)=n\operatorname{Re}z\) 在原点恒零，在别处发散。真正传播的是非负差函数的中心值，而不是任意函数的一个点值。


\ProseParagraphBreak

有界的整个调和函数也满足 Liouville 性质。设 \(|u|\le M\) 于平面，内部一阶导数估计在任意点、任意半径给 \(|\nabla u(a)|\le CM/R\)，令 \(R\to\infty\) 得梯度零，所以常数。实际上只要求单侧有界即可：若 \(u\ge0\)，对固定 \(z\) 在任意大圆盘应用 Harnack，不等式两端的系数均趋一，故 \(u(z)=u(0)\)。若只要求在某条直线有界，结论不成立，\(u(x+\mathrm{i}y)=y\) 已给反例。全平面的范围在扩大圆盘这一证明中被明确使用。

\ProseParagraphBreak

圆盘上的正调和函数还有边界测度表示，不要求其连续达到边界。对 \(0<s<R\)，把非负测度 \(u(se^{\mathrm{i}t})\mathrm{d}t/(2\pi)\) 放在参数圆周，它的总质量恒为 \(u(0)\)。紧圆周上的有限测度族具有弱收敛子列；此处可用连续函数可数稠密集上的对角抽取构造正泛函，再用正泛函表示得到极限测度 \(\nu\)。固定 \(|z|<R\)，在 \(s>|z|\) 的圆盘使用 Poisson 表示，并让子列 \(s\uparrow R\)，核关于参数一致收敛，得到
\[
u(z)=\int_0^{2\pi}P_R(z,t)\,\mathrm{d}\nu(t),\qquad \nu([0,2\pi))=u(0).
\]
傅里叶展开确定了 \(\nu\) 的全部三角矩，再由三角多项式稠密性确定测度，故表示唯一。这是边界数据由函数扩大到测度的一种自然方式。

\ProseParagraphBreak

特别地，取 \(\nu\) 为一点质量，就得到 \(P_R(z,t_0)\)。它在几乎所有边界点趋零，却在 \(Re^{\mathrm{i}t_0}\) 趋于无穷，内部仍严格正。因而“边界几乎处处为零”不等于经典连续边值为零；若不规定解的函数类，不能从几乎处处边值推出唯一性。Poisson 核的这个例子会在 Hardy 空间中再次出现，届时可积边界代表和测度奇异部分必须分别处理。


\ProseParagraphBreak

正调和函数的导数还可用自身点值控制。若 \(u\ge0\) 调和于 \(B(a,2r)\)，Harnack 给 \(\sup_{B(a,r)}u\le C u(a)\)，再在较小圆盘使用一阶内部估计，得到 \(|\nabla u(a)|\le C u(a)/r\)。当 \(u(a)>0\)，除以自身值可得 \(|\nabla\log u(a)|\le C/r\)。所以正调和函数在一个内部尺度上不能变化得任意快；常数只依赖尺度比而不依赖函数大小。若允许变号，\(u_n=x/n+y\) 等函数在某个点可有零值而梯度非零，便不可能由自身值给同样的界。

\ProseParagraphBreak

平均值与面积测度的归一化在缩放下保持一致。将 \(z=a+r\zeta\) 后，面积变为 \(r^2\mathrm{d}A(\zeta)\)，Laplace 算子变为 \(r^{-2}\Delta_\zeta\)，圆周长度变为 \(r\mathrm{d}s\)。因此圆周平均用 \(2\pi r\) 除，面积平均用 \(\pi r^2\) 除；分布基本解中的点质量则在对应变量替换时保持单位总质量。这些尺度检查会在后面的 Green 核、Riesz 测度与零点计数中反复出现。只看公式的函数部分而忽略测度变量，是丢失二维常数与半径幂的常见原因。
